% 埃尔德什·帕尔（Paul Erdős）（综述）
% license CCBYSA3
% type Wiki

本文根据 CC-BY-SA 协议转载翻译自维基百科\href{https://en.wikipedia.org/wiki/Paul_Erd\%C5\%91s}{相关文章}。

\begin{figure}[ht]
\centering
\includegraphics[width=6cm]{./figures/04eef634120f8175.png}
\caption{} \label{fig_ARPR_1}
\end{figure}
保罗·埃尔德什（英语：Paul Erdős，匈牙利语：Erdős Pál [ˈɛrdøːʃ ˈpaːl]；1913 年 3 月 26 日—1996 年 9 月 20 日）是一位匈牙利数学家。他是 20 世纪最具产出力的数学家之一，也是提出最多数学猜想的人之一。\(^\text{[2][3]}\)埃尔德什在离散数学、图论、数论、数学分析、逼近论、集合论以及概率论等领域追求并提出了大量问题。\(^\text{[4]}\)他的大部分研究集中在离散数学领域，解决了其中许多长期未解的难题。他大力倡导并作出重要贡献的拉姆齐理论，研究的是在何种条件下秩序必然出现。总体而言，埃尔德什的工作更倾向于解决已存在的开放问题，而非开辟或探索全新的数学领域。他一生共发表约 1500 篇数学论文，这一数字至今无人能及。\(^\text{[5]}\)

他以“社会化数学”的实践方式以及古怪的生活方式闻名；《时代》杂志称他为“怪人中的怪人”。\(^\text{[6]}\)他坚信数学是一项社会性活动，因此过着漂泊不定的生活，唯一目的就是与其他数学家共同撰写论文。即使在晚年，他仍将清醒的每个时刻都献给数学研究；1996 年，他在华沙参加数学会议时去世。\(^\text{[7]}\)

埃尔德什与合作者们的大量共同成果，促成了“埃尔德什数”的诞生——它表示某位数学家与埃尔德什之间的最短合著路径中的“步数”。
\subsection{生平}
保罗·埃尔德什于 1913 年 3 月 26 日出生在奥匈帝国的布达佩斯，\(^\text{[8]}\)是安娜（婚前姓威廉，Wilhelm）与拉约什·埃尔德什（原姓英兰德，Engländer）\(^\text{[9][10]}\)的独子。他的两个姐姐在他出生前几天因猩红热去世，年仅三岁和五岁。\(^\text{[11]}\)他的父母都是犹太人，任职于高中教授数学。埃尔德什从小便对数学表现出极大的兴趣。由于父亲在 1914 年至 1920 年间作为奥匈帝国的战俘被囚禁于西伯利亚，\(^\text{[10]}\)他幼年时期主要由一位德国女家庭教师照顾，\(^\text{[12]}\)母亲不得不长时间工作以维持生计。父亲在战俘营中自学了英语，但许多单词的发音不正确。后来父亲教儿子说英语时，保罗也继承了这种发音，并在其一生中都保持这种独特的口音。\(^\text{[13]}\)

他通过阅读父母放在家中的数学书籍自学阅读。五岁时，只要告诉他一个人的年龄，他就能在脑中计算出这个人一生到目前为止大约活了多少秒。\(^\text{[12]}\)由于两位姐姐早逝，他与母亲关系极为亲密，据说直到上大学前，他们一直同睡一张床。\(^\text{[14][15]}\)

16 岁时，父亲向他介绍了两个将成为他终生热爱的主题——无穷级数与集合论。在高中时期，埃尔德什热衷于解答《KöMaL》（《中学数学与物理杂志》）每月刊登的数学问题，并成为其中的活跃解题者之一。\(^\text{[16]}\)

埃尔德什在 17 岁时赢得全国考试后，进入布达佩斯大学学习。当时，匈牙利实行“入学限额制”（numerus clausus），对犹太学生的大学录取人数有严格限制。\(^\text{[13][17]}\)到他 20 岁时，他已找到伯特兰猜想（Bertrand’s postulate）的证明。\(^\text{[17]}\)1934 年，年仅 21 岁的他获得数学博士学位。\(^\text{[17]}\)他的博士导师是李波特·费耶尔（Lipót Fejér），这位著名数学家也是约翰·冯·诺伊曼、乔治·波利亚（George Pólya）以及保尔·图兰（Pál Turán）的导师。他随后在英国曼彻斯特大学获得博士后研究职位，因为当时匈牙利的犹太人正受到专制政权的压迫。在那里，他结识了 G.H. 哈代（G. H. Hardy）和斯坦尼斯瓦夫·乌拉姆（Stanisław Ulam）。\(^\text{[13]}\)

由于身为犹太人，匈牙利对他而言已不再安全，他于 1938 年离开祖国，移居美国。\(^\text{[17]}\)在第二次世界大战期间，埃尔德什的多位亲属在布达佩斯遇难，其中包括两位姑母、两位叔叔以及他的父亲，只有母亲幸存下来。当时，埃尔德什正在美国普林斯顿高等研究院（Institute for Advanced Study）任职。\(^\text{[17][18]}\)然而，由于他不符合该机构的行为标准（他们认为他“粗野且不合常规”\(^\text{[13]}\)），他的奖学金并未像预期那样延长一年，而只续期了六个月。

传记作者保罗·霍夫曼（Paul Hoffman）形容他“或许是世界上最古怪的数学家”。\(^\text{[19]}\)埃尔德什成年后的大部分人生都过着提着手提箱四处漂泊的生活。除了 1950 年代几年的特殊时期（他因被指“同情共产主义者”而被禁止进入美国），他几乎一生都在从一个会议或研讨会赶往另一个会议。\(^\text{[19]}\)每到一个地方，他都期望主人为他提供住宿、饮食、洗衣以及一切生活所需，并帮他安排前往下一站的行程。\(^\text{[19]}\)

1943 年，乌拉姆离开威斯康星大学麦迪逊分校，前往新墨西哥州洛斯阿拉莫斯参与“曼哈顿计划”，与其他数学家和物理学家共事。他曾邀请埃尔德什加入该项目，但当埃尔德什表达了战后想回匈牙利的意愿后，邀请被撤回。\(^\text{[13]}\)

1996 年 9 月 20 日，83 岁的埃尔德什在华沙参加学术会议时突发心脏病去世。\(^\text{[20]}\)他的离世几乎与他自己理想的方式一致。他曾说：

“我希望能在讲课时去世，当我在黑板上完成一个重要证明时，听众中有人喊道：‘那一般情形呢？’我会转身微笑着说：‘我把那留给下一代吧。’然后我就倒下。”\(^\text{[20]}\)

埃尔德什终身未婚，也没有子女。\(^\text{[9]}\)他安葬于布达佩斯的犹太人科兹马街公墓，与父母长眠于同一处。\(^\text{[21]}\)他曾为自己的墓志铭建议写上：“我终于不再变笨了。”（匈牙利语：“Végre nem butulok tovább”）。\(^\text{[22]}\)

埃尔德什的名字中包含匈牙利字母“ő”（双上尖音符的 o），但常被误写为 Erdos 或 Erdös，原因“要么是错误，要么是出于印刷上的需要”。\(^\text{[23]}\)
\subsection{职业生涯}
1934 年，埃尔德什前往英国曼彻斯特任客座讲师。1938 年，他接受了美国新泽西州普林斯顿高等研究院（Institute for Advanced Study）的奖学金职位，这是他在美国的第一个学术任职，并持续了接下来的十年。在此期间，他与马克·卡茨（Mark Kac）和奥雷尔·温特纳（Aurel Wintner）合作发表了关于概率数论的重要论文，与保尔·图兰（Pál Turán）合作研究逼近论，并与维托尔德·胡列维奇（Witold Hurewicz）合作研究维数理论。\(^\text{[24]}\)尽管成果卓著，他的奖学金并未续期，埃尔德什不得不以“流浪学者”的身份在宾夕法尼亚大学、圣母大学、普渡大学、斯坦福大学和雪城大学等多所高校间辗转任职。他从未在一个地方长期停留，而是在世界各地的数学机构之间不断旅行，直至去世。

由于当时的“红色恐慌”（Red Scare）与麦卡锡主义运动，\(^\text{[25][26][27]}\)1954 年，美国移民与归化局拒绝为埃尔德什（匈牙利公民）签发重入美国的签证。\(^\text{[28]}\)当时他正在圣母大学任教，本可选择留在美国，但他决定收拾行李离开，虽然此后仍定期向美国移民局提出重新审查的申请。此后，他一度移居以色列，在耶路撒冷希伯来大学获得为期三个月的任职，随后又在以色列理工学院（Technion）获得“永久访问教授”的职位。
\begin{figure}[ht]
\centering
\includegraphics[width=6cm]{./figures/a39475638355915e.png}
\caption{从左起逆时针方向：埃尔德什、范重（Fan Chung）及其丈夫罗纳德·格雷厄姆（Ronald Graham），摄于 1986 年日本。} \label{fig_ARPR_2}
\end{figure}
当时的匈牙利是华沙条约组织成员国，受苏联影响。虽然匈牙利政府严格限制本国公民的出入境自由，但在 1956 年，它破例授予埃尔德什特权——允许他自由地出入国境，不受限制。

1963 年，美国移民与归化局为埃尔德什签发了签证，他重新开始在美国的学术机构任教并往返访问。十年后的 1973 年，年满 60 岁的埃尔德什自愿离开匈牙利。\(^\text{[29]}\)

在生命的最后几十年中，埃尔德什共获得至少 15 个名誉博士学位，并成为包括美国国家科学院与英国皇家学会在内的 8 个国家科学院的成员。\(^\text{[30]}\)1977 年，他当选为荷兰皇家艺术与科学学院的外籍院士。\(^\text{[31]}\)去世前不久，他因认为同事阿德里安·邦迪（Adrian Bondy）受到了不公正对待，而主动放弃了加拿大滑铁卢大学授予他的名誉博士学位。\(^\text{[32][33]}\)
\subsubsection{数学研究}
埃尔德什是数学史上论文发表数量最多的学者之一，其产出量可与莱昂哈德·欧拉（Leonhard Euler）相媲美；不过两人风格不同——埃尔德什发表的论文更多（大多与他人合作完成），而欧拉撰写的页数更多（主要独立完成）。\(^\text{[34]}\)埃尔德什一生共发表约 1,525 篇数学论文，\(^\text{[35]}\)其中大部分为合作成果。他坚信数学是一种社会性活动，并以此为实践终身奉行，\(^\text{[36]}\)一生中共有 511 位不同的合作者。\(^\text{[37]}\)

在数学风格上，埃尔德什更倾向于“问题解决者”（problem solver）而非“理论构建者”（theory developer）。（有关这两种风格及其地位的详细讨论，可参见蒂莫西·高尔斯（Timothy Gowers）所著《数学的两种文化》（The Two Cultures of Mathematics）\(^\text{[38]}\)。）乔尔·斯宾塞（Joel Spencer）指出：“在 20 世纪数学史的殿堂中，埃尔德什的地位存在一定争议，因为他在辉煌的职业生涯中始终专注于具体定理和猜想。”\(^\text{[39]}\)埃尔德什从未获得菲尔兹奖（他在世时的最高数学荣誉），也从未与任何菲尔兹奖得主合作发表论文，\(^\text{[40]}\)这一规律同样延伸到其他奖项。\(^\text{[41]}\)不过，他在 1983/84 年获得了沃尔夫数学奖，以表彰“他在数论、组合学、概率论、集合论和数学分析方面的诸多贡献，以及他对全世界数学家的个人激励”。\(^\text{[42]}\)相比之下，获奖前后的三位得主的成果被分别描述为“杰出”“经典”“深刻”，以及“基础性”或“开创性”。

在埃尔德什的诸多贡献中，拉姆齐理论的发展以及概率方法的推广尤为突出。极值组合学（extremal combinatorics）作为一个重要分支，也在很大程度上受益于他的方法论，这一思路部分源自解析数论的传统。埃尔德什给出了伯特兰猜想（Bertrand’s postulate）的一个更为简洁的证明，比切比雪夫（Chebyshev）的原始证明更加优雅。他还与阿特勒·塞尔伯格（Atle Selberg）共同发现了素数定理（prime number theorem）的首个初等证明。然而，两人在证明过程与发表问题上的分歧最终导致了激烈的争执。\(^\text{[43][44]}\)此外，埃尔德什还在一些他本人兴趣不大的领域作出贡献，例如拓扑学——他首次构造出一个完全不连通但并非零维的拓扑空间，该空间后被称为“埃尔德什空间”（Erdős space）。\(^\text{[45]}\)
\subsubsection{埃尔德什的问题}
\begin{figure}[ht]
\centering
\includegraphics[width=6cm]{./figures/19f35a6cc5cac8ac.png}
\caption{埃尔德什对许多年轻数学家产生了深远影响。1985 年，在阿德莱德大学拍摄的这张照片中，埃尔德什正在向当时年仅 10 岁的陶哲轩讲解一道数学问题。陶哲轩于 2006 年获得菲尔兹奖，并于 2007 年当选为英国皇家学会院士。} \label{fig_ARPR_3}
\end{figure}
埃尔德什不仅以解决数学难题闻名，也以不断提出新问题而著称；数学家恩斯特·施特劳斯（Ernst Strauss）曾称他为“出题者中的绝对君主”。\(^\text{[7]}\)在其职业生涯中，埃尔德什经常为尚未解决的问题设立奖金。\(^\text{[46]}\)奖金额度从 25 美元（针对那些他认为稍微超出当时数学界理解范围的问题——包括他自己）到 10,000 美元不等，\(^\text{[47]}\)后者用于极其困难且具有重大数学意义的问题。其中一些问题后来被成功解决，包括奖金最高的那一道——“埃尔德什关于素数间隙的猜想”，于 2014 年被证明，\(^\text{[48]}\)奖金 10,000 美元也随之兑现。

据推测，至今仍有至少一千道“埃尔德什问题”尚未解决，尽管目前没有官方或完整的清单。即使在他去世后，这些悬赏依然有效；罗纳德·格雷厄姆（Ronald Graham）担任（非正式的）奖金管理员。解决者可以选择领取由埃尔德什生前签署、仅供纪念的原版支票（不可兑现），或由格雷厄姆签发的可兑现支票。\(^\text{[49]}\)[需要更新] 2024 年，英国数学家托马斯·布鲁姆（Thomas Bloom）创建了一个专门收录“埃尔德什问题”的网站。\(^\text{[50]}\)

其中最具数学意义的问题之一是埃尔德什关于算术级数的猜想（Erdős conjecture on arithmetic progressions）：

若一个整数序列的倒数和发散，则该序列必包含任意长度的算术级数。

如果此命题成立，它将解决数论中若干其他悬而未决的问题。其中一个主要推论——“素数集合中存在任意长度的算术级数”——后来已被独立证明，即格林–陶定理（Green–Tao theorem）。目前，该问题的悬赏金额为 5,000 美元。\(^\text{[51]}\)

最为大众熟知的“埃尔德什奖金问题”或许是考拉兹猜想（Collatz conjecture），也称“3N + 1 问题”。埃尔德什为该问题提供了 500 美元的奖金。\(^\text{[52]}\)
\subsubsection{合作者}
埃尔德什合作最频繁的伙伴包括匈牙利数学家安德拉什·沙尔柯齐（András Sárközy，合著 62 篇论文）和安德拉什·海纳尔（András Hajnal，合著 56 篇论文），以及美国数学家拉尔夫·福德里（Ralph Faudree，合著 50 篇论文）。其他经常合作的学者包括：\(^\text{[53]}\)
\begin{itemize}
\item 理查德·谢尔普（Richard Schelp，42 篇）
\item 塞西尔·C·鲁索（Cecil C. Rousseau，35 篇）
\item 维拉·T·绍什（Vera T. Sós，35 篇）
\item 阿尔弗雷德·雷尼（Alfréd Rényi，32 篇）
\item 保尔·图兰（Pál Turán，30 篇）
\item 恩德雷·塞梅雷迪（Endre Szemerédi，29 篇）
\item 罗恩·格雷厄姆（Ron Graham，28 篇）
\item 斯特凡·伯尔（Stefan Burr，27 篇）
\item 卡尔·波默朗斯（Carl Pomerance，23 篇）
\item 乔尔·斯宾塞（Joel Spencer，23 篇）
\item 亚诺什·帕赫（János Pach，21 篇）
\item 米克洛什·西蒙诺维茨（Miklós Simonovits，21 篇）
\item 恩斯特·G·施特劳斯（Ernst G. Straus，20 篇）
\item 梅尔文·B·纳森松（Melvyn B. Nathanson，19 篇）
\item 让-路易·尼古拉（Jean-Louis Nicolas，19 篇）
\item 理查德·拉多（Richard Rado，18 篇）
\item 贝拉·博洛巴什（Béla Bollobás，18 篇）
\item 埃里克·查尔斯·米尔纳（Eric Charles Milner，15 篇）
\item 安德拉什·杰尔法什（András Gyárfás，15 篇）
\item 约翰·塞尔弗里奇（John Selfridge，14 篇）
\item 范重（Fan Chung，14 篇）
\item 理查德·R·霍尔（Richard R. Hall，14 篇）
\item 乔治·皮拉尼安（George Piranian，14 篇）
\item 伊什特万·尤（István Joó，12 篇）
\item 佐尔特·图扎（Zsolt Tuza，12 篇）
\item A. R. 雷迪（A. R. Reddy，11 篇）
\item 沃伊泰赫·罗德尔（Vojtěch Rödl，11 篇）
\item 保尔·雷韦斯（Pál Révész，10 篇）
\item 佐尔坦·弗雷迪（Zoltán Füredi，10 篇）
\end{itemize}
关于其他与埃尔德什合作过的学者，可参见《埃尔德什数》条目中“埃尔德什数为 1 的人物名单”。
\subsection{埃尔德什数}
由于埃尔德什一生发表论文极为高产，他的朋友们为向他致敬而创造了“埃尔德什数”这一概念。埃尔德什数用来描述一个人与埃尔德什之间的“学术合作距离”——也就是他与埃尔德什本人，或与某位拥有自己埃尔德什数的合作者之间的合作层级。埃尔德什本人被定义为 0 号（即埃尔德什数为 0），与他直接合著论文的人拥有埃尔德什数 1，与这些人合作发表论文者则拥有最多 2 的埃尔德什数，以此类推。

目前，约有 20 万名数学家被分配了埃尔德什数，\(^\text{[54]}\)有研究估计，全球约 90\% 的活跃数学家其埃尔德什数小于 8——这一现象与“小世界效应”（small-world phenomenon）不谋而合。由于数学家与其他领域学者的跨学科合作，物理学、工程学、生物学与经济学等领域的许多科学家也拥有埃尔德什数。\(^\text{[55]}\)

多项研究显示，顶尖数学家往往拥有特别低的埃尔德什数。\(^\text{[56]}\)例如，在约 26.8 万名拥有埃尔德什数的数学家中，中位数为 5。\(^\text{[57]}\)相比之下，菲尔兹奖得主的中位数为 3。\(^\text{[58]}\)截至 2015 年，大约有 11,000 名数学家的埃尔德什数不超过 2。\(^\text{[59][60]}\)随着时间推移，合作距离将必然增长，因为拥有较低埃尔德什数的数学家逐渐去世，新的合作机会随之减少。美国数学学会（AMS）提供了一个免费的在线工具，可查询《数学评论》（Mathematical Reviews）目录中任意数学作者的埃尔德什数。\(^\text{[61]}\)

埃尔德什数最早可能由分析学家卡斯珀·戈夫曼（Casper Goffman）提出，\(^\text{[62]}\)他本人的埃尔德什数为 2——因为他曾与数学家理查德·B·达斯特（Richard B. Darst）合作，而达斯特又与埃尔德什合著论文。\(^\text{[63]}\)戈夫曼于 1969 年发表了一篇题为《你的埃尔德什数是多少？》（And what is your Erdős number?）的文章，对埃尔德什的高产合作现象进行了观察与分析。\(^\text{[64]}\)

数学家杰里·格罗斯曼（Jerry Grossman）曾幽默地指出，棒球名人堂成员汉克·阿伦（Hank Aaron）也可以被认为拥有 埃尔德什数 1，因为在埃默里大学授予他们名誉学位的同一天，他们共同在同一个棒球上签名送给了数学家卡尔·波默朗斯（Carl Pomerance）。\(^\text{[65]}\)此外，人们还为一名婴儿、一匹马以及若干演员提议计算“埃尔德什数”。\(^\text{[66]}\)
\subsection{性格}
物质财富对埃尔德什而言几乎毫无意义；由于他一生漂泊，他的所有财物几乎都能装进一个手提箱。他经常将获得的奖项奖金或其他收入捐赠给有需要的人或各种慈善事业。他的大部分人生都在世界各地的科学会议、大学以及同行的家中度过。作为客座讲师，他从各大学获得的津贴和奖项奖金足以支撑旅行与生活；剩余的钱，他则用来设立奖金，奖励解决“埃尔德什问题”的人（见上文）。

他常常突然出现在同事家门口，宣布一句：“我的大脑打开了（My brain is open）。”然后停留几天，与对方合作撰写几篇论文后便再次启程。在许多情况下，他会问合作伙伴：“接下来我该去找谁？”

他的同事阿尔弗雷德·雷尼（Alfréd Rényi）曾打趣道：“数学家是一台将咖啡转化为定理的机器。”\(^\text{[68]}\)埃尔德什确实酷爱咖啡，这句话常被误认为是他说的，\(^\text{[69]}\)但他自己将这句名言归功于雷尼。\(^\text{[70]}\)

1971 年母亲去世后，埃尔德什开始服用抗抑郁药与安非他命类兴奋剂，这令朋友们十分担忧。其中一位朋友罗纳德·格雷厄姆（Ron Graham）甚至与他打赌 500 美元，赌他能否一个月不碰药物。埃尔德什赢了这场赌局，但抱怨说戒药影响了他的思维：“你让我证明我不是瘾君子，但我一个月什么工作都没做。我每天早上起床对着空白纸发呆，脑中一片空白，简直像个普通人。你让数学倒退了整整一个月。”\(^\text{[71]}\)赢得赌注后，他立即恢复服用利他林（Ritalin）和苯丙胺（Benzedrine）。\(^\text{[72]}\)

埃尔德什有自己独特的语言体系。尽管他是不可知论的无神论者，\(^\text{[73][74]}\)却常提及一本象征性的“天书”（The Book），即上帝记录了所有数学定理最优美、最完美证明的那本书。\(^\text{[75]}\)他在 1985 年的一次讲座中说：“你不必相信上帝，但你应该相信《天书》。”他怀疑上帝的存在，\(^\text{[76][77]}\)还戏称上帝为“最高法西斯”（Supreme Fascist，简称 SF），责怪他“藏起了自己的袜子和匈牙利护照”，以及“把最优美的数学证明都据为己有”。当他看到特别优雅的数学证明时，会惊叹道：“这一定出自《天书》！”这一说法后来启发了著名的数学著作《来自天书的证明》（Proofs from the Book）。

埃尔德什独特的语言体系中还有许多其他富有个性的表达方式：\(^\text{[72]}\)
\begin{itemize}
\item 儿童被称为“ε（艾普西龙）”，因为在数学中（尤其是微积分中），任意小的正数通常用希腊字母 ε 表示。
\item 女性被称为“老板”（bosses），她们通过结婚“俘获”男性，使其成为“奴隶”；而离婚的男人则被称为“解放者”（liberated）。
\item 停止做数学的人被认为是“死了”，而真正去世的人则被称为“离开了”（left）。
\item 含酒精饮品被他称作“毒药”（poison）。
\item 音乐（除古典音乐外）都是“噪音”（noise）。
\item 被认为是平庸学者的人称为“牛顿”（Newton），暗讽缺乏创造性。
\item 发表数学讲座被称为“布道”（to preach）。
\item 数学讲座本身被称为“布道会”（sermons）。\(^\text{[78]}\)
\item 对学生进行口试则叫作“拷问”（to torture）他们。
\end{itemize}
他还给许多国家起了绰号，例如：美国被称为 “samland”（意指“山姆之国”，源自“山姆大叔” Uncle Sam）；苏联被称为 “joedom”（意指“约瑟夫之国”，源自约瑟夫·斯大林 Joseph Stalin）。\(^\text{[72]}\)他还打趣地说，印地语是“最好的语言”，因为“年老”（bud̩d̩hā）和“愚蠢”（buddhū）这两个词听起来几乎一样。\(^\text{[79]}\)
\subsubsection{签名}
埃尔德什在签名时常写作 “Paul Erdos P.G.O.M.”。当他年满 60 岁时，又在签名中添加了 “L.D.”；65 岁时加上 “A.D.”；70 岁时再次写上 “L.D.”；到 75 岁时则加上 “C.D.”。\(^\text{[79]}\)

这些缩写分别代表：
\begin{itemize}
\item P.G.O.M.：Poor Great Old Man（“贫穷而伟大的老人”）
\item L.D.（第一次）：Living Dead（“行尸走肉”）
\item A.D.：Archaeological Discovery（“考古发现”）
\item L.D.（第二次）：Legally Dead（“法律上已死”）
\item C.D.：Counts Dead（“统计学上的死亡”）\(^\text{[80][81]}\)
\end{itemize}
\subsection{遗产}
\subsubsection{书籍与电影}
\begin{figure}[ht]
\centering
\includegraphics[width=6cm]{./figures/e004d6c3e4e7cbf7.png}
\caption{埃尔德什之墓，布达佩斯科兹马街公墓} \label{fig_ARPR_4}
\end{figure}
埃尔德什至少成为了三本书的主角：两本传记——保罗·霍夫曼（Paul Hoffman）所著的《只爱数字的人》（The Man Who Loved Only Numbers）与布鲁斯·谢克特（Bruce Schechter）所著的《我的大脑是打开的》（My Brain is Open），这两部作品都于 1998 年出版；以及黛博拉·海利格曼（Deborah Heiligman）于 2013 年出版的儿童绘本《爱数学的男孩：保罗·埃尔德什的不凡人生》（The Boy Who Loved Math: The Improbable Life of Paul Erdős）。\(^\text{[82]}\)

他也是乔治·奇谢里（George Csicsery）拍摄的传记纪录片《数字是一个名字：保罗·埃尔德什肖像》（N is a Number: A Portrait of Paul Erdős）的主人公，该片在埃尔德什仍然在世时拍摄完成。\(^\text{[83]}\)
\subsubsection{天文学}
2021 年，小行星 405571 号（临时编号 2005 QE87）被正式命名为 “Erdőspál”，以纪念保罗·埃尔德什。命名说明中写道：“他是一位匈牙利数学家，其研究主要集中在离散数学领域。他的工作更倾向于解决现有的开放问题，而非开辟或探索新的数学领域。”\(^\text{[84]}\)该命名由 “K. Sárneczky” 与 “Z. Kuli”（后者为该小行星的发现者）共同提议。
\subsubsection{蛛形学}
2025 年，英国蛛形学家丹妮拉·舍伍德（Danniella Sherwood）与 R.C. 盖伦（R. C. Gallon）在尼日利亚发现并描述了一种新的狼蛛物种——尼氏异穴蛛（Heterothele erdosi Sherwood & Gallon, 2025），以此纪念保罗·埃尔德什。[85]
\subsection{另见}
\begin{itemize}
\item 以保罗·埃尔德什命名的主题列表——包括猜想、数字、奖项与定理
\item 造箱游戏（Box-making game）
\item 覆盖系统（Covering system）——同余类的集合
\item 图论维数（Dimension, graph theory）——与图相关的整数概念
\item 偶回路定理（Even circuit theorem）
\item 友谊图（Friendship graph）——由共享顶点的多个三角形组成的图
\item 匈牙利数学（Hungarian mathematics）——匈牙利数学的发展史
\item 最小重叠问题（Minimum overlap problem）
\item 概率方法（Probabilistic method）——用于数学证明的非构造性方法
\item 概率数论（Probabilistic number theory）——数论的一个分支
\item “火星人”科学家（The Martians, scientists）——指一群杰出的匈牙利科学家
\end{itemize}
\subsection{参考文献}
\begin{enumerate}
\item “Mathematics Genealogy Project.” 检索于 2012 年 8 月 13 日。
\item “The Sum-Product Problem Shows How Addition and Multiplication Constrain Each Other.” Quanta Magazine. 2019 年 2 月 6 日。检索于 2019 年 10 月 6 日。
\item Hoffman, Paul.（2013 年 7 月 8 日）。“Paul Erdős.” Encyclopædia Britannica.
\item “Paul Erdos - Hungarian mathematician.” Britannica.com. 检索于 2017 年 12 月 2 日。
\item 根据 “Facts about Erdös Numbers and the Collaboration Graph” 使用 *Mathematical Reviews 数据库的统计结果，论文发表数量排名第二的数学家约有 823 篇论文。
\item Lemonick, Michael D.（1999 年 3 月 29 日）。“Paul Erdos: The Oddball’s Oddball.” Time. 存档于 2012 年 1 月 6 日。
\item Kolata, Gina.（1996 年 9 月 24 日）。“Paul Erdos, 83, a Wayfarer In Math’s Vanguard, Is Dead.” The New York Times. 第 A1、B8 版。检索于 2008 年 9 月 29 日。
\item “Erdos biography.” Gap-system.org. 存档于 2011 年 6 月 7 日。检索于 2010 年 5 月 29 日。
\item Baker, A.; Bollobas, B.（1999）。 “Paul Erdős 26 March 1913 – 20 September 1996: Elected For.Mem.R.S. 1989.” Biographical Memoirs of Fellows of the Royal Society. 45: 147–164. doi:10.1098/rsbm.1999.0011。
\item Chern, Shiing-Shen; Hirzebruch, Friedrich.（2000）。Wolf Prize in Mathematics. World Scientific. 第 294 页。ISBN 978-981-02-3945-9。
\item “Paul Erdős.” 检索于 2015 年 6 月 11 日。
\item Hoffman 1998，第 66 页。
\item “Paul Erdős - Biography.” Maths History. 检索于 2022 年 7 月 6 日
\item Hoffman, Paul.（2016 年 7 月 1 日）。“Paul Erdős: The Man Who Loved Only Numbers” 视频讲座。YouTube. 曼彻斯特大学。检索于 2017 年 3 月 17 日。
\item Alexander, James.（1998 年 9 月 27 日）。“Planning an Infinite Stay.” The New York Times. 检索于 2022 年 5 月 6 日。
\item Babai, László. “Paul Erdős just left town.” 存档于 2011 年 6 月 9 日。
\item Bruno 2003，第 120 页。
\item Csicsery, George Paul.（2005）。N Is a Number: A Portrait of Paul Erdős. 柏林、海德堡：施普林格出版社。ISBN 3-540-22469-6。
\item Bruno 2003，第 121 页。
\item Bruno 2003，第 122 页。
\item “Erdős Pál sírja - grave 17A-6-29.” agt.bme.hu.存档于 2016 年 4 月 4 日。检索于 2017 年 12 月 2 日。
\item Hoffman 1998，第 3 页。
\item 引文原文为：“请注意字母 ‘ő’ 上的两个长音符号，它常常（甚至在埃尔德什自己的论文中）因错误或印刷需要被替换为 ‘ö’，即更常见的德语变音符，而这种符号在匈牙利语中也存在。” 出自：Erdős, Paul; Miklós, D.; Sós, Vera T. (1996). Combinatorics, Paul Erdős is eighty.
\item Bollobás 1996，第 4 页。
\item “The wandering mathematician: Paul Erdos.” TheArticle.*2023 年 7 月 28 日。检索于 2023 年 9 月 9 日。
\item “Paul Erdős - Biography.” Maths History. 检索于 2023 年 9 月 9 日。
\item Sack, Harald.（2018 年 9 月 20 日）。“What’s your Erdös Number? – The bustling Life of Mathematician Paul Erdös.” SciHi Blog. 存档于 2023 年 11 月 23 日。检索于 2023 年 9 月 10 日。
\item “Erdos biography.” *School of Mathematics and Statistics, University of St Andrews, Scotland.* 2000 年 1 月。检索于 2008 年 11 月 11 日。
\item Babai, László; Spencer, Joel. “Paul Erdős (1913–1996)” (PDF). *Notices of the American Mathematical Society.* 第 45 卷，第 1 期。美国数学学会。
\item Baker, A.; Bollobás, B.（1999）。 “Paul Erdős. 26 March 1913 — 20 September 1996.” *Biographical Memoirs of Fellows of the Royal Society.* 第 45 卷。伦敦皇家学会：147–164。doi:10.1098/rsbm.1999.0011。ISSN 0080-4606。S2CID 123517792。
\item “P. Erdös (1913 - 1996).” *Royal Netherlands Academy of Arts and Sciences.* 存档于 2020 年 7 月 28 日。
\item Erdős, Paul.（1996 年 6 月 4 日）。“Dear President Downey.” (PDF). 存档于 2005 年 10 月 15 日。检索于 2014 年 7 月 8 日。\\
“我怀着沉重的心情决定与滑铁卢大学断绝关系，包括辞去我在 1981 年获得的名誉学位（那一荣誉曾让我非常高兴）。我对艾德里安·邦迪教授（Adrian Bondy）受到的待遇感到极度不安。我并不声称邦迪教授完全无辜，但考虑到他对大学的成就与卓越贡献，我认为‘正义应当带着宽容’。”
\item 1996 年 10 月 2 日滑铁卢大学校报 University of Waterloo Gazette 文章转录（存档于 2010 年 11 月 23 日，Wayback Machine）。
\item Hoffman 1998，第 42 页。
\item Grossman, Jerry. “Publications of Paul Erdös.” 检索于 2011 年 2 月 1 日。
\item Krauthammer, Charles.（1996 年 9 月 27 日）。“Paul Erdos, Sweet Genius.” The Washington Post. 第 A25 页。
\item “The Erdős Number Project Data Files.”Oakland.edu. 2009 年 5 月 29 日。检索于 2010 年 5 月 29 日。
\item Gowers, Timothy.（2000）。 “The Two Cultures of Mathematics.” (PDF). 收录于 Arnold, V. I.; Atiyah, Michael; Lax, Peter D.; Mazur, Barry（编），Mathematics: Frontiers and Perspectives. 美国数学学会。ISBN 978-0821826973。
\item Spencer, Joel.（2000 年 11–12 月）。“Prove and Conjecture!” American Scientist.第 88 卷第 6 期。此文为《Mathematics: Frontiers and Perspectives》一书的书评。
\item  “Paths to Erdős - The Erdős Number Project - Oakland University.” oakland.edu. 检索于 2017 年 12 月 2 日。
\item 出自 The Mathematical Intelligencer第 21 卷第 3 期（1999 年夏季），第 51–63 页，作者 DeCastro 与 Grossman，《Trails to Erdos》（存档于 2015 年 9 月 24 日 Wayback Machine）。\\
文中表 3 的仔细分析表明，尽管埃尔德什从未与 42 位菲尔兹奖得主合作发表论文（这一点或许值得深思），但还有许多其他重要的国际数学奖项，例如罗尔夫·涅万林纳奖（Rolf Nevanlinna Prize）、沃尔夫数学奖（Wolf Prize in Mathematics）以及勒罗伊·斯蒂尔奖（Leroy P. Steele Prizes）。……也许令人好奇的是，卡普兰斯基（Kaplansky）是这些奖项获得者中唯一一位与埃尔德什合作过的数学家。（该论文发表后，合作者洛瓦兹（Lovász）获得了沃尔夫奖，使得人数增至两位。）
\item “Wolf Foundation Mathematics Prize Page.” Wolffund.org.il.存档于 2008 年 4 月 10 日。检索于 2010 年 5 月 29 日。
\item Goldfeld, Dorian（2003）。〈素数定理的初等证明：历史视角〉（The Elementary Proof of the Prime Number Theorem: an Historical Perspective）。载于 Number Theory: New York Seminar，第 179–192 页。
\item Baas, Nils A.; Skau, Christian F.（2008）。〈数字之主：阿特勒·塞尔伯格——他的生平与数学〉（The Lord of the Numbers, Atle Selberg. On his Life and Mathematics）（PDF）。Bulletin of the American Mathematical Society，第 45 卷第 4 期，617–649 页。doi:10.1090/S0273-0979-08-01223-8。
\item Henriksen, Melvin.〈对保罗·埃尔德什（1913–1996）的回忆〉（Reminiscences of Paul Erdös (1913–1996)）。美国数学协会（Mathematical Association of America）。检索于 2008 年 9 月 1 日。
\item “Math genius left unclaimed sum.” Edmonton Journal. 存档于 2011 年 1 月 18 日。检索于 2020 年 7 月 16 日。
\item “Prime Gap Grows After Decades-Long Lull.” 2014 年 12 月 10 日。
\item Kevin Hartnett（2017 年 6 月 5 日）。〈Cash for Math: The Erdős Prizes Live On〉。
\item Seife, Charles（2002 年 4 月 5 日）。〈Erdös’s Hard-to-Win Prizes Still Draw Bounty Hunters〉。Science. 第 296 卷第 5565 期，39–40 页。doi:10.1126/science.296.5565.39。PMID 11935003。S2CID 34952867。
\item “Erdős Problems.” Erdős Problems. 检索于 2024 年 4 月 23 日。
\item Soifer, Alexander（2008）。The Mathematical Coloring Book: Mathematics of Coloring and the Colorful Life of its Creators. 纽约：施普林格出版社，第 354 页。ISBN 978-0-387-74640-1。
\item Kurtz, Stuart A.; Simon, Janos（2007）。〈广义考拉兹问题的不可判定性〉（The Undecidability of the Generalized Collatz Problem）。载于 Cai, Jin-Yi; Cooper, S. Barry; Zhu, Hong（编），Theory and Applications of Models of Computation. Lecture Notes in Computer Science, 第 4484 卷。柏林、海德堡：施普林格出版社，第 542–553 页。doi:10.1007/978-3-540-72504-6_49。ISBN 978-3-540-72504-6。
\item “List of collaborators of Erdős by number of joint papers.” The Erdős Number Project 网站。存档于 2008 年 8 月 4 日（Wayback Machine）。
\item “From Benford to Erdös.” Radio Lab. 节目日期 2009 年 10 月 9 日。2009 年 9 月 30 日播出。存档于 2010 年 8 月 18 日。检索于 2016 年 2 月 6 日。
\item Grossman, Jerry. “Some Famous People with Finite Erdös Numbers.” 检索于 2011 年 2 月 1 日。
\item De Castro, Rodrigo; Grossman, Jerrold W.（1999）。〈通往保罗·埃尔德什的著名路径〉（Famous Trails to Paul Erdős）（PDF）。The Mathematical Intelligencer. 第 21 卷第 3 期，51–63 页。CiteSeerX 10.1.1.33.6972。doi:10.1007/BF03025416。MR 1709679。S2CID 120046886。原文存档于 2015 年 9 月 24 日。西班牙文原版发表于 Rev. Acad. Colombiana Cienc. Exact. Fís. Natur. 第 23 卷第 89 期（1999），第 563–582 页，MR 1744115。
\item “Facts about Erdös Numbers and the Collaboration Graph.” The Erdös Number Project – Oakland University. OU 主页面。检索于 2019 年 10 月 6 日。

\end{enumerate}